\documentclass[10pt,a4paper]{amsart}
\usepackage[margin=19mm]{geometry}
\usepackage{amsmath,amssymb,mathtools}
\usepackage[scr=rsfs,cal=euler]{mathalfa}
\usepackage{xfrac}
\usepackage[unicode,hidelinks,bookmarks=false]{hyperref}
\newcommand{\ie}{i.e.\ }
\newcommand{\cf}{cf.\ }
\newcommand{\resp}{respectively\ }
\newcommand{\Subsection}{\subsection}
\providecommand{\nobreakdash}{\nobreak\hbox{-}}
\setlength{\emergencystretch}{2em}
\multlinegap0pt
\usepackage{amssymb}
\usepackage[shortlabels]{enumitem}
\usepackage[normalem]{ulem}
\usepackage{tikz}
\newtheorem{assumption}{Assumption}
\newtheorem{proposition}{Proposition}
\newcommand{\C}{\mathbb{C}}
\newcommand{\LP}{\mathcal{LP}}
\newcommand{\R}{\mathbb{R}}
\renewcommand\Re{\operatorname{Re}}
\newcommand{\dd}{\mathrm{d}}
\DeclareMathOperator{\ii}{i}
\DeclareMathOperator{\supp}{supp}
\title{P\'olya--Schur problems and free probability}
\author{Andrew Campbell, Jonas Jalowy}
\date{Source: arXiv:2605.31356v1 (CC BY 4.0)}
\begin{document}
\maketitle

\noindent\emph{Source and licence.} P\'olya--Schur problems and free probability. Version arXiv:2605.31356v1 (CC BY 4.0), distributed by arXiv under the Creative Commons Attribution 4.0 International licence, http://creativecommons.org/licenses/by/4.0/, as recorded in the arXiv licence metadata for this exact version and shown on the abstract page https://arxiv.org/abs/2605.31356v1. Full licence notice: \texttt{licenses/arxiv-2605.31356-CC-BY-4.0.txt}.\par\smallskip

\noindent\textit{Editorial scope.} Counted unit: the article's proposition R-transforms (source0.tex lines 714--732). The article's complete original proof of it is delivered with it (source0.tex lines 816--841, one \texttt{proof} environment of 26 lines, which the article places away from the statement). No mathematics is written, repaired or re-derived: the statement and the proof are the article's own lines, in the article's own notation, and no retained line is substituted or re-flowed.  Retained uncounted as the source-local context the proof needs: lines 310--320.  Nothing was omitted from any retained span, so the package carries no omission record.  No retained line contains a citation, so the wrapper carries no bibliography and no bibliography entry.  No retained line contains a figure environment, an image-inclusion command or drawing code, so no image is delivered and the package declares no asset dependency.  Editorial formatting only, in addition: an AMS \texttt{amsart} preamble that restates the article's own declarations and macros used by the retained lines, namely the environments \texttt{assumption}, \texttt{proposition} on separate counters, together with the standard packages those lines need.  The retained environments print in this document as Assumption 1, Proposition 1 in order of appearance, whereas the article prints the same items with the numbering of its own full document.  The complete native source of the article as distributed in the arXiv e-print for this exact version is bundled unchanged as \texttt{sources/201716/source0.tex}.
\begin{assumption} \label{assump:main assumption}
	Let $0\leq M_n=o(n)$, and let $\{f_{n}\}_{n=1}^{\infty}$ be such that $f_n\in \LP_{M_n}$ and of the form \begin{equation}
		f_n(z)=e^{-c_{n}z-\frac{\sigma_{n}^2}{2}z^2}\prod_{j=1}^{\infty}\left(1-\alpha_{j,n}z\right)e^{\alpha_{j,n}z}.
	\end{equation} 
	We say the sequence $\{f_n\}$ satisfies Assumption \ref{assump:main assumption} if the finite complex valued measures
	\begin{equation}\label{eq:Sigma_n assumptions}
		\Sigma_{n}:=n\sigma_{n}^2\delta_{0}+\frac{1}{n}\sum_{j=1}^{\infty}\frac{n^2\alpha_{j,n}^2}{n^2\alpha_{j,n}^2+1}\delta_{n\alpha_{j,n}} \to \Sigma
	\end{equation}
	converge vaguely to some positive weighted L\'evy measure $\Sigma$ as $n\to\infty$, and if 
	\begin{equation}\label{eq:c assumptions} \lim\limits_{n\rightarrow\infty}c_{n}-\int t\dd\Sigma_{n}(t)=c.\end{equation} %{eq:f_n assumptions}
	\end{assumption}

 \begin{proposition}\label{prop:convergence of R-transforms}
	Let $\{f_{n}\}_{n\geq 1}$ satisfy Assumption \ref{assump:main assumption}. 
	\begin{enumerate}
		\item Let $R_{\mu_{c,\Sigma}}$ be the $R$-transform of $\mu_{c,\Sigma}$, then 
		\begin{equation}
			\lim\limits_{n\rightarrow\infty}-\frac{f_{n}'(nu)}{f_{n}(nu)}= R_{\mu_{c,\Sigma}}(u),
		\end{equation} 
		uniformly on compact subsets of $\C_{-}$.
		
		\item Define $\left[ \log(1-xw)\frac{x^2+1}{x^2}+\frac{w}{x}\right]\Big|_{x=0}:=-\frac{w^2}{2}$ by continuous extension, then		
		\begin{equation}\label{eq:log of f_n convergence}
			\lim\limits_{n\rightarrow\infty}\frac{1}{n}\log f_{n}(nu)-\frac{1}{n}\log f_{n}(-n\ii)=-c(u+\ii)+\int_{\R}\log\left(\frac{1-ux}{1+\ii x}\right)\frac{x^2+1}{x^2}+\frac{u+\ii}{x}\dd\Sigma(x),
		\end{equation}  uniformly on compact subsets of $\C_{-}$.
		\item For any $M\in\R$ define $\log_{M}|z|:=\max(\log|z|,-M)$, then \begin{align}\label{eq:cutoff log of f_n convergence}
			\lim\limits_{n\rightarrow\infty} &\frac{1}{n}\log_{M}| f_{n}(nu)|-\frac{1}{n}\log_{M}| f_{n}(-n\ii)|\\
			&=\Re\left(-c(u+\ii)\right)+\int_{\R}\left[\log_{M}\left|{1-ux}\right|-\log_{M}\left|{1+\ii x}\right|\right]\frac{x^2+1}{x^2} +\Re\frac{u+\ii}{x}\dd\Sigma(x),\nonumber
		\end{align} uniformly on compact subsets of $\C_{-}\cup\R\setminus\{0\}$.
	\end{enumerate} 
\end{proposition} 

\begin{proof}[Proof of Proposition \ref{prop:convergence of R-transforms}]
	For any $u\in\C_{-}$ \begin{equation}
		\begin{aligned}
			-\frac{f_{n}'(nu)}{f_{n}(nu)}&=c_{n}+n\sigma_{n}^2u+\sum_{j=1}^{\infty}\frac{\alpha_{j,n}}{1-n\alpha_{j,n}u}-\alpha_{j,n}\\
			&=c_n+n\sigma_{n}^{2}u+\frac{1}{n}\sum_{j=1}^{\infty}\frac{n\alpha_{j,n}}{1-n\alpha_{j,n}u}-n\alpha_{j,n}\\
			&=c_n-\int_{\C}x\dd\Sigma_{n}(x)+\int_{\C}\frac{u+x}{1-xu}\dd\Sigma_{n}(x).
		\end{aligned}
	\end{equation}  In general, $\supp\left(\Sigma_{n}\right)=\{0\}\cup\{n\alpha_{j,n}\}\subset\C$. Note, by assumption  $f_{n}\in\LP_{M_n}$ for some $0\leq M_n=o(n)$, and we can strengthen this bound on the support to  \begin{equation}\label{eq:supp_inclusion}
		\supp\left(\Sigma_{n}\right)=\{0\}\cup\{n\alpha_{j,n}\}\subset \C\setminus \left(B_{\frac{n}{2M_n} }\left(\frac{-\ii n}{2M_n}\right)\cup B_{\frac{n}{2M_n} }\left(\frac{\ii n}{2M_n}\right) \right),
	\end{equation} where $B_{r}(w)$ is an open disk of radius $r$ centered at $w$. Note that the right hand side of \eqref{eq:supp_inclusion} is precisely the image of $\R+\ii[-M_n/n,M_n/n]$ under the inversion map. Eventually $\frac{1}{u}$ is contained in these open disks away from the support of $\Sigma_{n}$. Thus, $h_{u}(x)=\frac{u+x}{1-xu}$ is a bounded continuous function on $\bigcup_{n\geq N}\supp\left(\Sigma_{n}\right)$ for some $N$, and moreover the bound is uniform for $u$ in compact subsets of $\C_{-}$ with the subspace topology. Thus, by \eqref{eq:Sigma_n assumptions} of Assumption \ref{assump:main assumption} \begin{equation}
		\begin{aligned}
			\lim\limits_{n\rightarrow\infty}-\frac{f_{n}'(nu)}{f_{n}(nu)}&= c+ \int_{\R}\frac{u+x}{1-xu} \dd\Sigma(x), %c+\sigma^2u+\int_{\R}\frac{u+t}{1-tu}\frac{t^2}{t^2+1}\dd\nu(t),\\
		\end{aligned}
	\end{equation} which is the $R$-transform of the $\mu_{c,\Sigma}$. Moreover, this convergence is uniform on compact subsets of $\C_{-}$ by Montel's Theorem.
	
	The proof of \eqref{eq:log of f_n convergence} follows in a similar manner after recalling we define $$\left[ \log(1-xw)\frac{x^2+1}{x^2}+\frac{w}{x}\right]\Big|_{x=0}:=-\frac{w^2}{2}.$$ Specifically, \begin{equation}
		\begin{aligned}
			\frac{1}{n}\log f_n(nu)-\frac{1}{n}\log f_{n}(-n\ii)&=-\frac{n\sigma_n^2}{2}u^2-c_nu+\frac{1}{n}\sum_{j=1}^{\infty}\log\left(1-un\alpha_{j,n}\right)+un\alpha_{j,n}\\
			&\quad+\frac{n\sigma_n^2}{2}(-\ii)^2+c_n\ii-\frac{1}{n}\sum_{j=1}^{\infty}\log\left(1+\ii n\alpha_{j,n}\right)-\ii n\alpha_{j,n}\\
			&=-(u+\ii)\left[c_n-\int x\dd\Sigma_{n}(x)\right]\\
			&\quad +\int_{\C}\left[\log(1-ux)-\log(1+\ii x)\right]\frac{x^2+1}{t^2}+\frac{u+\ii}{x}\dd\Sigma_{n}(x)
		\end{aligned}
	\end{equation} The integrand is continuous away from $x=\frac{1}{u}$ and by the same argument as above is  a continuous bounded function on some domain containing $\bigcup_{n\geq N}\supp\Sigma_{n}$ for some large $N$. \eqref{eq:log of f_n convergence} then follows from Assumption \ref{assump:main assumption}.
	
	Finally, \eqref{eq:cutoff log of f_n convergence} follows in a nearly identical manner to \eqref{eq:log of f_n convergence} after noting that the integrand is continuous and bounded for any $u\in\C_-\cup\R\setminus\{0\}$, and this bound is uniform on compact subsets of $\C_-\cup\R\setminus\{0\}$. 
\end{proof}

\end{document}
